\begin{textsample}{POS Dim 3 – vem\_ed\_07 – Score 7.00 – vem\_ed\_07/t026.txt}  \label{ex:f3_pos_076}
Brasil e Colômbia em sintonia Em 2020, após uma série de reuniões entre a equipe dos PCIs da Cesu / Centro Paula Souza e da universidade \textbf{colombiana} Uniminuto, começou uma fértil colaboração entre as duas instituições.
No segundo semestre de 2020, o professor Joelson Alves do Nascimento (Fatec Barueri) realizou um intercâmbio virtual com o colega Luigi Ramirez (Uniminuto) sobre competências emocionais.
Após Instagram, LinkedIn e o quebra-gelo, com troca de informações sobre a \textbf{faculdade}, sobre lazer e cultura, 23 estudantes de Gestão da Tecnologia da Informação compartilharam com os alunos \textbf{colombianos} artigos em espanhol e português sobre o tema das competências emocionais no ambiente de trabalho: empatia, comprometimento, assertividade, relacionamento interpessoal e comunicação.
As parcerias prosseguem em 2021: no primeiro semestre, Lilian de Souza (Fatec Piracicaba) e Miguel Angel Nicholls (Uniminuto) abordaram a troca de informações entre culturas.
As equipes de estudantes pesquisaram temas como festas, danças, comidas, lazer e educação.
Mais sobre o projeto em https://sites.google.com/view/eciaespanol/.
Regiane Moreira (Fatec Guaratinguetá) e Mónica Paola Diaz Oliveros (Uniminuto) desenvolveram o PCI"Influência do jargão e uso dos recursos linguísticos e semióticos na publicidade digital", com 24 estudantes brasileiros e 8 \textbf{colombianos}.
As equipes \textbf{mistas} estudaram peças publicitárias nas plataformas Facebook, YouTube, Instagram, LinkedIn e Glassdoor, analisando o uso da língua e seus efeitos nas relações sociais e interpessoais, os aspectos semióticos visuais e seus significados possíveis dentro das culturas, encontrando semelhanças e diferenças a partir do visual nos dois países.
Elizabeth Herrera Colorado e Tânia Augusta Ferreira (Fatec Itaquaquecetuba) e Mayra Alejandra Rodríguez (Uniminuto) trabalharam"Competência linguística e Marketing"com 26 estudantes no Brasil e 21 na Colômbia.
Os objetivos do projeto foram fazer contato com outras culturas, conhecer o \textbf{consumidor} e o mercado internacional, aplicar conceitos de \textbf{marketing} a possíveis situações de novos negócios e planejar lançamento de um produto brasileiro na Colômbia e vice-versa.
As colaborações também incluíram o Colóquio de Literatura Brasileira: olhares e \textbf{impressões}, realizado em 14 de maio de 2021 pelas Fatecs Guaratinguetá e Itaquaquecetuba para estudantes da universidade \textbf{colombiana} (leia mais na página 3).

% matched lemmas: colombiano, consumidor, faculdade, impressão, marketing, mista
\end{textsample}
